% ============================================
% LUC-04
% Gerado de molambudos.md — build_miolo.py
% ============================================
\section*{\textbf{LUC-04} --- A Fotografia}
\markboth{LUC-04 --- A Fotografia}{}

\noindent \textit{O velho no cemitério}

\noindent \textbf{Março de 2026 --- Barbacena}

\vspace{0.5em}\noindent\rule{\textwidth}{0.5pt}\vspace{0.5em}

\noindent Lúcia voltou ao hotel às 22h.
\noindent O ar entrou grosso. Por um segundo, ela teve a impressão absurda de que não respirava: era respirada.

\noindent O hotel era uma pousada simples no centro de Barbacena. Quarto pequeno, cama dura, uma janela que dava para a rua principal. Ela colocou a mochila na cama, tirou o diário e o caderneta, e sentou.

\noindent O diário estava sobre a cama.

\noindent Ela ainda não tinha aberto.

\noindent Nas duas semanas desde que pegou o diário no Arquivo Público, ela não tinha conseguido abrir. Sempre havia uma desculpa: trabalho, cansaço, hora errada. A verdade era que ela estava com medo.

\noindent Não medo racional.
Medo do que ia encontrar.
Medo do que ia sentir.
\noindent Não medo racional. Medo do que ia encontrar. Medo do que ia sentir. Medo do que ia mudar dentro dela.

\noindent \textit{\textquotedblleft{}Você é psiquiatra forense\textquotedblright{}\textit{, ela disse em voz alta. }\textquotedblleft{}Você lida com criminosos, com traumatizados, com arquivos de horror. Isso é só mais um caso.\textquotedblright{}}

\noindent Ela abriu o diário.

\noindent A primeira página tinha a caligrafia infantil de Joaquim. A data: 1917.

\noindent \textit{\textquotedblleft{}Hoje eu cheguei no Colônia. O médico disse que eu sou doido. Mas eu não sou doido. A criatura é que é doida. Ela mora dentro de mim e fala o tempo todo.\textquotedblright{}}

\noindent Lúcia leu.

\noindent Página após página.

\noindent A infância no sertão. A seca. A caminhada. O pai morto no juazeiro. A mãe levada pelos soldados. O Curral do Governo. A vala.

\noindent Ela leu sobre a criatura.

\noindent Leu a descrição do corpo de espinhos, das costelas viradas para fora, do órgão pulsando. Leu sobre a luz presa no olho. Leu sobre o trato.

\noindent Leu sobre o olho amarelo. Leu sobre o trato.

\noindent Quando levantou os olhos, eram 3:14 da manhã.

\noindent O quarto estava escuro. A única luz vinha do abajur da mesa de cabeceira. O diário estava aberto na página 126.

\noindent Ela sentiu o cheiro.

\noindent Adocicado.

\noindent Estava no quarto.

\noindent Ela olhou para o canto. A sombra parecia mais densa ali. Mais escura. Como se o escuro fosse mais escuro naquele pedaço da parede.

\noindent \textit{\textquotedblleft{}Não tem nada\textquotedblright{}}, ela disse.

\noindent Mas não acreditou.

\noindent Ela fechou o diário. Guardou na mochila. Deitou. Apagou a luz.

\noindent Ficou olhando para o teto.

\noindent Às 3:14 da manhã, o olho amarelo apareceu no canto do quarto.

\noindent Não era um sonho.

\noindent Ela viu.

\noindent Piscou.

\noindent O olho sumiu.

\noindent Ela passou o resto da noite com a luz acesa, o diário fechado debaixo do travesseiro, e a certeza de que alguma coisa tinha entrado no quarto.

\noindent A certeza de que alguma coisa tinha entrado \textit{nela}.

\vspace{0.5em}\noindent\rule{\textwidth}{0.5pt}\vspace{0.5em}

\noindent Na manhã seguinte, ela voltou ao Colônia.

\noindent Dessa vez, foi até o jazigo. O cemitério do hospital, onde sessenta mil corpos foram enterrados em valas comuns. Não havia cruzes. Não havia nomes. Apenas um gramado liso, como se nada estivesse enterrado ali.

\noindent Não havia cruzes. Não havia nomes. Apenas um gramado liso, como se nada estivesse enterrado ali.

\noindent Ela sabia que Joaquim estava ali. Cova 1.260.

\noindent Ela ficou parada, olhando para o chão.

\noindent \textit{\textquotedblleft{}O que você quer de mim?\textquotedblright{}}, perguntou em voz alta.

\noindent O vento respondeu com silêncio.

\noindent \textit{\textquotedblleft{}Por que eu?\textquotedblright{}}

\noindent Uma folha seca rodopiou aos pés dela.

\begin{figure}[H]
  \centering
  \includegraphics[width=0.55\textwidth]{foto_05_fotografia_luc04_sepia.png}
  \caption*{\footnotesize\textbf{Fotografia 4} --- \textit{O velho no cemitério. Jazigo do Colônia, 2026.}}
\end{figure}

\noindent \textit{\textquotedblleft{}Você leu\textquotedblright{}}, disse uma voz atrás dela.

\noindent Lúcia se virou.

\noindent Um homem velho estava parado a alguns metros. Muito velho. Pele enrugada, cabelos brancos ralos, olhos opacos. Ele segurava uma bengala.

\noindent \textit{\textquotedblleft{}Quem é você?\textquotedblright{}}

\noindent \textit{\textquotedblleft{}Só um velho que passou por aqui\textquotedblright{}\textit{, ele disse. }\textquotedblleft{}E que viu a senhora lendo o diário.\textquotedblright{}}

\noindent \textit{\textquotedblleft{}Como sabe que eu li?\textquotedblright{}}

\noindent Ele sorriu. A boca era uma linha torta.

\noindent \textit{\textquotedblleft{}O cheiro. A senhora está com o cheiro. Quem leu o diário fica com o cheiro por dias.\textquotedblright{}}

\noindent \textit{\textquotedblleft{}Que cheiro?\textquotedblright{}}

\noindent \textit{\textquotedblleft{}Doce. De flor murcha. De carne passada.\textquotedblright{}}

\noindent Lúcia sentiu o estômago apertar.

\noindent \textit{\textquotedblleft{}Quem é você?\textquotedblright{}} repetiu.

\noindent O velho se aproximou. Devagar. A bengala batendo no chão seco.

\noindent \textit{\textquotedblleft{}Conheci o paciente 1.260.\textquotedblright{}\textit{ }\textquotedblleft{}Mentira. Ele morreu em 79.\textquotedblright{}}

\noindent \textit{\textquotedblleft{}Conheci antes. Trabalhei aqui nos anos 50.\textquotedblright{}}

\noindent \textit{\textquotedblleft{}O que você sabe sobre ele?\textquotedblright{}}

\noindent O velho parou na frente dela.

\noindent \textit{\textquotedblleft{}Sei que ele escrevia num diário. Sei que o diário foi enterrado com ele. Sei que, quando exumaram, o diário estava quente. E sei que todo mundo que leu o diário...\textquotedblright{}\textit{ Ele fez uma pausa. }\textquotedblleft{}...teve visita.\textquotedblright{}}

\noindent \textit{\textquotedblleft{}Que visita?\textquotedblright{}}

\noindent \textit{\textquotedblleft{}A Rasga Mortalha.\textquotedblright{}}

\noindent O velho apontou para o céu.

\noindent Lúcia olhou.

\noindent No galho seco de uma árvore, uma coruja branca estava pousada. Ela olhava fixamente para Lúcia.

\noindent \textquotedblleft{\textit{Ela vem quando alguém abre o diário\textquotedblright{}}, disse o velho. }\textquotedblleft{\textit{Fica olhando até a pessoa escolher.\textquotedblright{}}}

\noindent \textit{\textquotedblleft{}Que ciclo?\textquotedblright{}}

\noindent \textit{\textquotedblleft{}O ciclo do paciente 1.260. Ele morre, alguém lê, o leitor vira o próximo.\textquotedblright{}}

\noindent O velho deu as costas.

\noindent \textquotedblleft{\textit{Espera!\textquotedblright{}} Lúcia chamou. }\textquotedblleft{\textit{Como eu paro isso?\textquotedblright{}}}

\noindent Ele parou. Não virou o rosto.

\noindent \textit{\textquotedblleft{}Não sai.\textquotedblright{}}

\noindent \textit{\textquotedblleft{}Não sai?\textquotedblright{}}

\noindent \textit{\textquotedblleft{}Só passa adiante. Arranja outro leitor. Passa o diário. Deixa outra pessoa ler. Aí o ciclo continua. Só que com outro.\textquotedblright{}}

\noindent Ele começou a andar.

\noindent \textit{\textquotedblleft{}E se eu não passar?\textquotedblright{}}

\noindent A voz do velho veio de longe, carregada pelo vento:
\noindent \textquotedblleft{\textit{Não para. Troca de mão. A senhora fecha a sua e outra abre. É assim que coisa enterrada aprende a andar.\textquotedblright{}}}

\noindent \textquotedblleft{\textit{Então a senhora apodrece com ele por dentro. Igual Joaquim. Igual o médico. Igual todos que acharam que podiam guardar uma porta fechada no peito.\textquotedblright{}}}

\noindent A coruja branca abriu as asas. Voou.

\noindent Lúcia ficou sozinha no cemitério, com o diário na mochila, o cheiro no nariz, e a certeza de que alguma coisa havia mudado para sempre.

\vspace{0.5em}
\noindent\textit{\textbf{↪ Rotas:} \rota{LUC-05} \rotasep \rota{MEM-12}}